% !TEX spellcheck = en_US



% ~~~~~~~~~~~~~~~~~~~~~~~~~~~~~~~~~~~~~~~ V
% (NC) 2026 Haluk Bingol
% github.com/halukbingol/LaTeX-Templates
%
% Licensees may copy, distribute, display, and perform the work and make derivative works 
% and remixes based on it only for non-commercial purposes. 
% ~~~~~~~~~~~~~~~~~~~~~~~~~~~~~~~~~~~~~~~ A




% =====================================================================
%  Java on the Command Line -- Senior Level: Packaging, Start-up, Diagnostics and Supply Chain
%  ---------------------------------------------------------------------
%  Built on hbCisLaTeXTemplate.tex with the libraries in ../../hblib/.
%  Programs and their outputs are read from codeoutput/:
%     codeoutput/<Name>.java   the program
%     codeoutput/<Name>.in     keyboard input (optional)
%     codeoutput/<Name>.txt    what it printed, made by:  sh run_examples.sh
% =====================================================================

\documentclass[aspectratio=169,10pt,compress]{beamer}

% ~~~~~~~~~~~~~~~~~~~~~~~~~~~~~~~~~~~~~~~~~~~~~~~~~~~~~~~~~~~~~~~~~ V hblib
	\usepackage{../../hblib/hblibPresentation} % lib presentation: theme, i/N, \hSection, algorithm2e, ...
	\usepackage{../../hblib/hblibCodeoutput} % lib code + output: \lstcode, \lstcodedown, ...
	\usepackage{../../hblib/hblibDatastructures} % lib data structures: \hString, \hStack, \hQueue
% ~~~~~~~~~~~~~~~~~~~~~~~~~~~~~~~~~~~~~~~~~~~~~~~~~~~~~~~~~~~~~~~~~ A hblib


% xcolor, array (and the L{width} column), listings, tikz and algorithm2e come from the hblib libraries
\usepackage[T1]{fontenc}
\usepackage{lmodern}
\usepackage{booktabs}

\definecolor{gitred}{RGB}{240,80,50}
\newcommand{\cmd}[1]{\texttt{\color{gitred}#1}}

\tikzset{
  box/.style={draw, thick, rounded corners, align=center, font=\small, inner sep=5pt},
  tool/.style={-{Stealth[length=2.5mm]}, thick}
}

% ---------------------------------------------------------------------
\title{Java on the Command Line}
\subtitle{Senior Level: Packaging, Start-up, Diagnostics and Supply Chain}
\author[Bingol]{Haluk O. Bingol}

\institute[FCIS]{
    Faculty of Computer and Information Sciences\\
    Yeditepe University
}
\date{
	{\tiny \hbTimeStamp}
}

\begin{document}




% ~~~~~~~~~~~~~~~~~~~~~~~~~~~~~~~~~~~~~~ V frm
\begin{frame}[t,fragile]
  \titlepage
\end{frame}
% ~~~~~~~~~~~~~~~~~~~~~~~~~~~~~~~~~~~~~~ A frm




% ~~~~~~~~~~~~~~~~~~~~~~~~~~~~~~~~~~~~~~ V frm
\begin{frame}[t,fragile] \frametitle{Outline}
  \tableofcontents[hideallsubsections]
\end{frame}
% ~~~~~~~~~~~~~~~~~~~~~~~~~~~~~~~~~~~~~~ A frm

% ~~~~~~~~~~~~~~~~~~~~~~~~~~~~~~~~~~~~~~ V
\hSection[Ship]{Shipping an Application}{}{}
% ~~~~~~~~~~~~~~~~~~~~~~~~~~~~~~~~~~~~~~ A

% ~~~~~~~~~~~~~~~~~~~~~~~~~~~~~~~~~~~~~~ V
\subsection[Options]{Options}
% ~~~~~~~~~~~~~~~~~~~~~~~~~~~~~~~~~~~~~~ A




% ~~~~~~~~~~~~~~~~~~~~~~~~~~~~~~~~~~~~~~ V frm
\begin{frame}[t,fragile] \frametitle{From \texttt{.class} files to something users can run}
  \begin{center}
  \begin{tikzpicture}[x=1cm, y=1cm]
    \node[box, fill=blue!8] (c) at (0,0) {classes\\ \scriptsize\texttt{out/}};
    \node[box, fill=yellow!20] (j) at (3,0) {jar\\ \scriptsize\texttt{app.jar}};
    \node[box, fill=orange!12] (r) at (6.4,0) {runtime image\\ \scriptsize\texttt{jlink}};
    \node[box, fill=green!12] (p) at (10.2,0) {app image / installer\\ \scriptsize\texttt{jpackage}};
    \draw[tool] (c) -- node[above, font=\ttfamily\scriptsize]{jar} (j);
    \draw[tool] (j) -- node[above, font=\ttfamily\scriptsize]{jdeps+jlink} (r);
    \draw[tool] (r) -- (p);
  \end{tikzpicture}
  \end{center}
  \small
  \begin{tabular}{@{}L{0.2\textwidth}L{0.3\textwidth}L{0.42\textwidth}@{}}
    \toprule
    \textbf{Form}          & \textbf{User needs}                   & \textbf{Typical use} \\
    \midrule
    executable jar         & a matching Java installed             & servers, tools for developers, coursework \\
    \texttt{jlink} image   & nothing; one OS and CPU per image     & containers, internal tools \\
    \texttt{jpackage}      & nothing; double-click or installer    & desktop apps: \texttt{.msi}/\texttt{.exe}, \texttt{.dmg}/\texttt{.pkg}, \texttt{.deb}/\texttt{.rpm} \\
    container image        & a container runtime                   & cloud services (\texttt{FROM eclipse-temurin:21-jre}) \\
    \bottomrule
  \end{tabular}
\end{frame}
% ~~~~~~~~~~~~~~~~~~~~~~~~~~~~~~~~~~~~~~ A frm

% ~~~~~~~~~~~~~~~~~~~~~~~~~~~~~~~~~~~~~~ V
\subsection[jpackage]{jpackage}
% ~~~~~~~~~~~~~~~~~~~~~~~~~~~~~~~~~~~~~~ A




% ~~~~~~~~~~~~~~~~~~~~~~~~~~~~~~~~~~~~~~ V frm
\begin{frame}[t,fragile] \frametitle{\texttt{jpackage}: a self-contained application}
  \lstoutput[basicstyle=\ttfamily\tiny]{codeoutput/Jpackage.txt}

  \vspace{0.2em}
  \small Default image: a full runtime (157\,MB). With a \texttt{jlink} runtime of only \texttt{java.base}: 39\,MB.
  The app image runs without an installed Java. Recorded on Linux; the next slide shows Windows and macOS.
\end{frame}
% ~~~~~~~~~~~~~~~~~~~~~~~~~~~~~~~~~~~~~~ A frm




% ~~~~~~~~~~~~~~~~~~~~~~~~~~~~~~~~~~~~~~ V frm
\begin{frame}[t,fragile] \frametitle{\texttt{jpackage} on Windows and macOS}
  \begin{columns}[T]
    \begin{column}{0.48\textwidth}
      \begin{block}{Windows}
        \footnotesize


% +++++++++++++++++++++++++++++++++++++++ V lst
%: -Lst.
\begin{lstlisting}[style=codesnippet, language={}, basicstyle=\ttfamily\tiny]
jpackage --type msi --name Hello ^
  --input input --main-jar app.jar ^
  --win-menu --win-shortcut
\end{lstlisting}
% +++++++++++++++++++++++++++++++++++++++ A lst



        \begin{itemize}
          \item \texttt{msi} and \texttt{exe} need the \textbf{WiX Toolset} on the build machine
          \item \texttt{\^} continues a line in Command Prompt
          \item Unsigned installers trigger SmartScreen; sign with \texttt{signtool}
        \end{itemize}
      \end{block}
    \end{column}
    \begin{column}{0.48\textwidth}
      \begin{block}{macOS}
        \footnotesize


% +++++++++++++++++++++++++++++++++++++++ V lst
%: -Lst.
\begin{lstlisting}[style=codesnippet, language={}, basicstyle=\ttfamily\tiny]
jpackage --type dmg --name Hello \
  --input input --main-jar app.jar \
  --mac-sign --mac-signing-key-user-name "..."
\end{lstlisting}
% +++++++++++++++++++++++++++++++++++++++ A lst



        \begin{itemize}
          \item \texttt{dmg} or \texttt{pkg}; Xcode command line tools needed
          \item Gatekeeper blocks unsigned apps; sign and \textbf{notarize} for distribution
          \item Build on Apple silicon for \emph{aarch64}, on Intel for \emph{x64}
        \end{itemize}
      \end{block}
    \end{column}
  \end{columns}
  \small \hRemark{No cross-building}: \texttt{jpackage} makes packages only for the operating system it runs on (use CI runners per OS).
\end{frame}
% ~~~~~~~~~~~~~~~~~~~~~~~~~~~~~~~~~~~~~~ A frm

% ~~~~~~~~~~~~~~~~~~~~~~~~~~~~~~~~~~~~~~ V
\hSection[Startup]{Start-up Time}{}{}
% ~~~~~~~~~~~~~~~~~~~~~~~~~~~~~~~~~~~~~~ A

% ~~~~~~~~~~~~~~~~~~~~~~~~~~~~~~~~~~~~~~ V
\subsection[CDS]{Class data sharing}
% ~~~~~~~~~~~~~~~~~~~~~~~~~~~~~~~~~~~~~~ A




% ~~~~~~~~~~~~~~~~~~~~~~~~~~~~~~~~~~~~~~ V frm
\begin{frame}[t,fragile] \frametitle{Where start-up time goes}
  \begin{columns}[T]
    \begin{column}{0.5\textwidth}
      \small Before \texttt{main} runs, the JVM must
      \begin{enumerate}\small
        \item start (memory, GC, threads)
        \item \textbf{load, verify and link} hundreds of classes
        \item run static initializers; the code is \textbf{interpreted} until the JIT compiles hot methods
      \end{enumerate}
      \textbf{CDS} (class data sharing) stores step 2's result in an archive that is memory-mapped at start.
      The JDK ships one for its own classes; \textbf{AppCDS} adds your application's classes.
    \end{column}
    \begin{column}{0.46\textwidth}
      \begin{block}{Two runs (Java 13+)}
        \footnotesize
        \texttt{java -XX:ArchiveClassesAtExit=app.jsa -jar app.jar}\\
        \texttt{java -XX:SharedArchiveFile=app.jsa -jar app.jar}\\[3pt]
        or one flag (Java 19+): \texttt{-XX:+AutoCreateSharedArchive -XX:SharedArchiveFile=app.jsa}
      \end{block}
      \small Further: Project Leyden (ahead-of-time class loading and linking, JDK 24+), GraalVM \texttt{native-image}
      (compiles to a native executable; closed-world limits on reflection).
    \end{column}
  \end{columns}
\end{frame}
% ~~~~~~~~~~~~~~~~~~~~~~~~~~~~~~~~~~~~~~ A frm




% ~~~~~~~~~~~~~~~~~~~~~~~~~~~~~~~~~~~~~~ V frm
\begin{frame}[t,fragile] \frametitle{Measuring CDS and AppCDS}
  \begin{columns}[T]
    \begin{column}{0.56\textwidth}
      \lstcode[basicstyle=\ttfamily\tiny]{codeoutput/proj/CDS/src/app/Startup.java}{Java}
    \end{column}
    \begin{column}{0.42\textwidth}
      \uncover<2->{\lstoutput[basicstyle=\ttfamily\tiny]{codeoutput/CDS.txt}}
    \end{column}
  \end{columns}

  \vspace{0.2em}
  \small Over 800 classes come from the JDK's CDS archive; without it start-up nearly doubles.
  AppCDS helps most with thousands of classes (frameworks). One run each; PowerShell: \texttt{Measure-Command}.
\end{frame}
% ~~~~~~~~~~~~~~~~~~~~~~~~~~~~~~~~~~~~~~ A frm

% ~~~~~~~~~~~~~~~~~~~~~~~~~~~~~~~~~~~~~~ V
\hSection[Diagnostics]{Diagnosing a Running JVM}{}{}
% ~~~~~~~~~~~~~~~~~~~~~~~~~~~~~~~~~~~~~~ A

% ~~~~~~~~~~~~~~~~~~~~~~~~~~~~~~~~~~~~~~ V
\subsection[jcmd]{jps, jcmd, JFR}
% ~~~~~~~~~~~~~~~~~~~~~~~~~~~~~~~~~~~~~~ A




% ~~~~~~~~~~~~~~~~~~~~~~~~~~~~~~~~~~~~~~ V frm
\begin{frame}[t,fragile] \frametitle{Tools for a running JVM}
  \small
  \begin{tabular}{@{}L{0.32\textwidth}L{0.6\textwidth}@{}}
    \toprule
    \textbf{Command}                                & \textbf{Shows} \\
    \midrule
    \texttt{jps -l}                                 & Java processes of this user (pid, main class) \\
    \texttt{jcmd <pid|class> help}                  & all diagnostic commands of that JVM \\
    \texttt{jcmd ... Thread.print}                  & stack of every thread (like \texttt{jstack}): hangs, deadlocks \\
    \texttt{jcmd ... GC.heap\_info}, \texttt{GC.class\_histogram} & heap use, objects per class: memory leaks \\
    \texttt{jcmd ... GC.heap\_dump file.hprof}      & full heap dump for a profiler (Eclipse MAT, VisualVM) \\
    \texttt{jcmd ... JFR.start}, \texttt{JFR.dump}   & start/save a flight recording in production \\
    \texttt{jfr summary/print rec.jfr}              & read a recording; or open it in JDK Mission Control \\
    \bottomrule
  \end{tabular}

  \vspace{0.5em}
  \small Same commands on Windows and macOS. They attach only to JVMs of the \textbf{same user}; in containers run them inside the container.
\end{frame}
% ~~~~~~~~~~~~~~~~~~~~~~~~~~~~~~~~~~~~~~ A frm




% ~~~~~~~~~~~~~~~~~~~~~~~~~~~~~~~~~~~~~~ V frm
\begin{frame}[t,fragile] \frametitle{Inspecting a busy program}
  \begin{columns}[T]
    \begin{column}{0.34\textwidth}
      \lstcode[basicstyle=\ttfamily\tiny]{codeoutput/proj/Diagnose/Busy.java}{Java}
    \end{column}
    \begin{column}{0.64\textwidth}
      \uncover<2->{\lstoutput[basicstyle=\ttfamily\tiny]{codeoutput/Diagnose.txt}}
    \end{column}
  \end{columns}

  \vspace{0.2em}
  \small \texttt{jcmd} accepts the main class instead of a pid. \texttt{jfr view} (Java 21) summarizes the recording:
  every CPU sample is in \texttt{Busy.work}. JFR costs about 1\,\% overhead, so it is used in production.
\end{frame}
% ~~~~~~~~~~~~~~~~~~~~~~~~~~~~~~~~~~~~~~ A frm

% ~~~~~~~~~~~~~~~~~~~~~~~~~~~~~~~~~~~~~~ V
\subsection[Hooks]{Shutting down}
% ~~~~~~~~~~~~~~~~~~~~~~~~~~~~~~~~~~~~~~ A




% ~~~~~~~~~~~~~~~~~~~~~~~~~~~~~~~~~~~~~~ V frm
\begin{frame}[t,fragile] \frametitle{Stopping cleanly: signals and shutdown hooks}
  \lstcode[basicstyle=\ttfamily\scriptsize]{codeoutput/proj/Hooks/Server.java}{Java}
  \uncover<2->{\lstoutput[basicstyle=\ttfamily\scriptsize]{codeoutput/Hooks.txt}}

  \vspace{0.2em}
  \small \texttt{kill -TERM} (and \texttt{Ctrl}+\texttt{C}, \texttt{docker stop}) runs the shutdown hooks; exit code $143 = 128 + 15$ (SIGTERM).
  \texttt{kill -9} skips them. On Windows, \texttt{Ctrl}+\texttt{C} runs the hooks, \texttt{taskkill /F} does not.
\end{frame}
% ~~~~~~~~~~~~~~~~~~~~~~~~~~~~~~~~~~~~~~ A frm

% ~~~~~~~~~~~~~~~~~~~~~~~~~~~~~~~~~~~~~~ V
\hSection[Supply]{Supply Chain and Builds}{}{}
% ~~~~~~~~~~~~~~~~~~~~~~~~~~~~~~~~~~~~~~ A

% ~~~~~~~~~~~~~~~~~~~~~~~~~~~~~~~~~~~~~~ V
\subsection[Sign]{Signing jars}
% ~~~~~~~~~~~~~~~~~~~~~~~~~~~~~~~~~~~~~~ A




% ~~~~~~~~~~~~~~~~~~~~~~~~~~~~~~~~~~~~~~ V frm
\begin{frame}[t,fragile] \frametitle{Signing and verifying a jar}
  \lstoutput[basicstyle=\ttfamily\tiny]{codeoutput/Sign.txt}

  \vspace{0.2em}
  \small The signature covers every entry's digest. After tampering, verification fails. Real releases use a CA-issued
  certificate and a \textbf{timestamp} (\texttt{-tsa}) so signatures stay valid after the certificate expires.
\end{frame}
% ~~~~~~~~~~~~~~~~~~~~~~~~~~~~~~~~~~~~~~ A frm

% ~~~~~~~~~~~~~~~~~~~~~~~~~~~~~~~~~~~~~~ V
\subsection[Repro]{Reproducible builds}
% ~~~~~~~~~~~~~~~~~~~~~~~~~~~~~~~~~~~~~~ A




% ~~~~~~~~~~~~~~~~~~~~~~~~~~~~~~~~~~~~~~ V frm
\begin{frame}[t,fragile] \frametitle{Reproducible jars (\texttt{jar --date}, since JDK 17.0.3)}
  \lstoutput[basicstyle=\ttfamily\tiny]{codeoutput/Repro.txt}

  \vspace{0.2em}
  \small Jars store file times, so two builds of the same code differ (\texttt{a}, \texttt{b}). With a fixed \texttt{--date}
  the bytes are identical (\texttt{c}, \texttt{d}): anyone can check that a published jar came from the published source.
  macOS: \texttt{shasum -a 256}; Windows: \texttt{certutil -hashfile c.jar SHA256}.
\end{frame}
% ~~~~~~~~~~~~~~~~~~~~~~~~~~~~~~~~~~~~~~ A frm




% ~~~~~~~~~~~~~~~~~~~~~~~~~~~~~~~~~~~~~~ V frm
\begin{frame}[t,fragile] \frametitle{Supply-chain checklist}
  \begin{columns}[T]
    \begin{column}{0.48\textwidth}
      \begin{block}{Build}
        \small
        \begin{itemize}
          \item Pin the JDK version (\texttt{--release}, CI image, toolchains)
          \item Pin dependency versions; verify checksums or signatures
          \item Reproducible output (\texttt{jar --date}, Maven \texttt{project.build.outputTimestamp})
          \item Generate an SBOM (software bill of materials, CycloneDX)
        \end{itemize}
      \end{block}
    \end{column}
    \begin{column}{0.48\textwidth}
      \begin{block}{Run}
        \small
        \begin{itemize}
          \item Smallest runtime that works (\texttt{jlink}): fewer modules, smaller attack surface
          \item Keep the JDK patched (quarterly updates for LTS releases)
          \item Do not run as root/Administrator
          \item Secrets via environment or files, never as \texttt{-D} on the command line (visible in \texttt{ps}, \texttt{jcmd VM.system\_properties})
        \end{itemize}
      \end{block}
    \end{column}
  \end{columns}
\end{frame}
% ~~~~~~~~~~~~~~~~~~~~~~~~~~~~~~~~~~~~~~ A frm

% ~~~~~~~~~~~~~~~~~~~~~~~~~~~~~~~~~~~~~~ V
\hSection[Containers]{Java in Containers}{}{}
% ~~~~~~~~~~~~~~~~~~~~~~~~~~~~~~~~~~~~~~ A

% ~~~~~~~~~~~~~~~~~~~~~~~~~~~~~~~~~~~~~~ V
\subsection[Memory]{Memory and CPUs}
% ~~~~~~~~~~~~~~~~~~~~~~~~~~~~~~~~~~~~~~ A




% ~~~~~~~~~~~~~~~~~~~~~~~~~~~~~~~~~~~~~~ V frm
\begin{frame}[t,fragile] \frametitle{The JVM reads container limits}
  \lstoutput[basicstyle=\ttfamily\tiny]{codeoutput/Container.txt}

  \vspace{0.2em}
  \small The JVM detects cgroup memory and CPU limits (\texttt{UseContainerSupport}, on by default). The default heap is only
  25\,\% of the limit; services often use \texttt{-XX:MaxRAMPercentage=75}. Docker Desktop on Windows and macOS runs a Linux VM,
  so these flags behave the same there.
\end{frame}
% ~~~~~~~~~~~~~~~~~~~~~~~~~~~~~~~~~~~~~~ A frm

% ~~~~~~~~~~~~~~~~~~~~~~~~~~~~~~~~~~~~~~ V
\hSection[Summary]{Summary}{}{}
% ~~~~~~~~~~~~~~~~~~~~~~~~~~~~~~~~~~~~~~ A

% ~~~~~~~~~~~~~~~~~~~~~~~~~~~~~~~~~~~~~~ V
\subsection[Cheat sheet]{Cheat sheet}
% ~~~~~~~~~~~~~~~~~~~~~~~~~~~~~~~~~~~~~~ A




% ~~~~~~~~~~~~~~~~~~~~~~~~~~~~~~~~~~~~~~ V frm
\begin{frame}[t,fragile] \frametitle{Cheat sheet}
  \footnotesize
  \begin{tabular}{@{}L{0.56\textwidth}L{0.36\textwidth}@{}}
    \toprule
    \textbf{Command} & \textbf{Does} \\
    \midrule
    \texttt{jpackage --type app-image|msi|dmg|deb ...}            & self-contained app or installer \\
    \texttt{jpackage ... --runtime-image rt}                       & use a small \texttt{jlink} runtime \\
    \texttt{java -XX:ArchiveClassesAtExit=a.jsa ...}               & create an AppCDS archive \\
    \texttt{java -XX:SharedArchiveFile=a.jsa ...}                  & start faster with it \\
    \texttt{jps -l}, \texttt{jcmd <pid> Thread.print}              & find a JVM, dump its threads \\
    \texttt{java -XX:StartFlightRecording=filename=r.jfr ...}      & record; \texttt{jfr summary r.jfr} \\
    \texttt{keytool -genkeypair ...}, \texttt{jarsigner}           & sign and verify jars \\
    \texttt{jar --date=2026-01-01T00:00:00Z ...}                   & reproducible jar \\
    \texttt{java -XX:MaxRAMPercentage=75 ...}                      & heap relative to container memory \\
    \bottomrule
  \end{tabular}
\end{frame}
% ~~~~~~~~~~~~~~~~~~~~~~~~~~~~~~~~~~~~~~ A frm

% ~~~~~~~~~~~~~~~~~~~~~~~~~~~~~~~~~~~~~~ V
\subsection[Projects]{Projects}
% ~~~~~~~~~~~~~~~~~~~~~~~~~~~~~~~~~~~~~~ A




% ~~~~~~~~~~~~~~~~~~~~~~~~~~~~~~~~~~~~~~ V frm
\begin{frame}[t,fragile] \frametitle{Projects}
  \begin{enumerate}\small
    \item Ship a course project three ways: executable jar, \texttt{jlink} image, \texttt{jpackage} installer (on your OS); compare size and start-up.
    \item Write one CI workflow (e.g.\ GitHub Actions) that builds \texttt{.msi} on a Windows runner and \texttt{.dmg} on a macOS runner.
    \item Measure start-up of a larger program (e.g.\ with many classes) with and without AppCDS; repeat 10 times and report mean and spread.
    \item Introduce a deadlock between two threads; find it with \texttt{jcmd <pid> Thread.print}.
    \item Make your build reproducible end to end; prove it by building on Windows and macOS and comparing SHA-256 hashes.
    \item Run a service in a container with \texttt{--memory=256m}; find the heap the JVM chose, then tune it.
  \end{enumerate}
\end{frame}
% ~~~~~~~~~~~~~~~~~~~~~~~~~~~~~~~~~~~~~~ A frm




% ~~~~~~~~~~~~~~~~~~~~~~~~~~~~~~~~~~~~~~ V frm
\begin{frame}[t,fragile]
	\begin{center}
		\vspace{4cm}
		\hRemark{\Huge Questions?}
	\end{center}
\end{frame}
% ~~~~~~~~~~~~~~~~~~~~~~~~~~~~~~~~~~~~~~ A frm



\end{document}
